% !TEX root = ../main.tex
% ============================================================================
% Chapter 5: 对比分析与讨论
% 数据来源: outputs/figures/10~13_*.png
% ============================================================================

\section{对比分析与讨论}

前两章分别评估了传统机器学习与神经网络方法在 NSL-KDD 数据集上的表现，
但各章实验在测试划分、调优策略、随机种子等方面存在差异，模型之间的横向比较难以直接进行。
本章将所有模型置于同一评测框架下重新评估，量化分析不同方法族之间的真实差距，
并结合训练耗时、可解释性、部署成本等维度给出综合的选型建议。

\subsection{评估方法}

为保证横向对比的公平性，所有模型均在统一测试集（22,544 条样本，20 维特征）上重新推理与指标计算。
测试集来自 NSL-KDD 官方 KDDTest+.txt，包含 15 种训练阶段从未出现的新型攻击，是评估模型开放世界泛化能力的客观基准。

二分类任务统一报告 accuracy、precision、recall、f1、AUC 五项指标；
多分类任务则采用 accuracy 与 f1\_macro（宏平均 F1，对类别不均衡敏感）两项指标，各类别 per-class F1 详见附录表7。
已有研究在 NSL-KDD 上进行了类似的对比实验~\cite{udurume2024comparative,candan2025ensemble,sunil2023comparative}，本文在此基础上进一步统一测试集、消除验证集偏差，确保各类方法在同一评测基准下公平对比。

\subsection{二分类性能对比}

表5 汇总了 6 个模型在统一测试集上的二分类指标，
其中 DT 与 RF 为第 3 章网格搜索调优后的最优模型，
MLP 同时给出基线与调优两组结果，CNN 与 LSTM 为第 4 章固定架构的最终模型。

\begin{table*}[t]
  \centering
  \caption{六模型二分类统一测试集指标对比}
  \label{tab:binary-comparison}
  \small
  \begin{tabular}{lccccc}
    \toprule
    模型 & Accuracy & Precision & Recall & F1 & AUC \\
    \midrule
    DT (tuned)  & \textbf{0.7755} & 0.9604 & \textbf{0.6317} & \textbf{0.7621} & 0.7821 \\
    RF (tuned)  & 0.7234 & 0.9602 & 0.5363 & 0.6882 & \textbf{0.9150} \\
    MLP (base)  & 0.7423 & 0.9580 & 0.5724 & 0.7166 & 0.9058 \\
    MLP (tuned) & 0.7445 & 0.9560 & 0.5777 & 0.7202 & 0.9074 \\
    CNN         & 0.7635 & \textbf{0.9635} & 0.6075 & 0.7452 & 0.8820 \\
    LSTM        & 0.7427 & 0.9628 & 0.5700 & 0.7161 & 0.8832 \\
    \bottomrule
  \end{tabular}
\end{table*}

对比呈现两个明显特征。第一，DT 调优后取得最高 accuracy（0.7755）和最高 F1（0.7621），
但其 AUC 仅 0.7821，是六个模型中最低的。这一"准确率高、判别力低"的悖论源于 DT 的硬决策特性：
树模型一旦做出预测即输出 0 或 1，缺乏概率校准，使得 ROC 曲线在高阈值区域快速塌陷。
第二，RF 的 F1 仅 0.6882，但 AUC 高达 0.9150，是所有模型中概率排序能力最强的。
RF 由 300 棵子树集成的投票机制天然输出校准良好的后验概率，
在阈值敏感的业务场景（如风险评级、欺诈拦截）中具有不可替代的优势。

这一对比揭示了精确率-召回率权衡（precision-recall tradeoff）的本质：
高精度（>0.95）几乎成为所有模型的共识，意味着误报控制良好；
真正拉开差距的是召回率（0.54–0.63），即"在保证低误报的前提下能多检出多少真异常"。
DT 以稍高的召回率换取了更高的 F1，RF 则以牺牲召回换取最佳的全局判别能力。

图14 以并列柱状图形式直观呈现传统机器学习与深度学习方法
在 accuracy、precision、recall、F1、AUC 五个维度上的差异，可视化地印证了上述分析。

\begin{figure*}
  \centering
  \includegraphics[width=\textwidth]{13_dl_vs_ml_comparison.png}
  \caption{传统机器学习与深度学习方法多指标对比}
  \label{fig:13_dl_vs_ml_comparison}
\end{figure*}

\subsection{ROC曲线分析}

ROC 曲线通过变化分类阈值描绘真阳性率与假阳性率的权衡，是评估模型概率排序能力的标准工具。
图15 将六个模型的 ROC 曲线叠加在同一坐标系中，
曲线下面积（AUC）直接反映了模型的全局判别能力。

\begin{figure*}
  \centering
  \includegraphics[width=\textwidth]{11_roc_curves.png}
  \caption{六模型 ROC 曲线叠加对比}
  \label{fig:11_roc_curves}
\end{figure*}

AUC 由高到低排序为：RF (0.9150) $>$ MLP-tuned (0.9074) $>$ MLP-base (0.9058) $>$ LSTM (0.8832) $>$ CNN (0.8820) $>$ DT (0.7821)。
曲线呈现两个聚类：RF 与两个 MLP 变体在低 FPR 区域紧贴左上角，AUC 均在 0.90 以上，
代表"概率质量集中于真实正例"的优秀排序；
CNN 与 LSTM 紧随其后，AUC 在 0.88 附近；
DT 因硬决策输出概率集中在 0 与 1 两个端点，ROC 曲线呈现明显的阶梯状，
导致曲线下面积被显著低估。这进一步验证了第 3 章的观察：
DT 适合作为高精度、低推理成本的实时检测器，但不适合作为概率评分模型。

\subsection{攻击大类F1对比}

为深入对比多分类模型在各类别上的细粒度表现，图16 给出了
DT、RF、MLP 三个多分类模型在 5 大类（DoS、Normal、Probe、R2L、U2R）上的 per-class F1 柱状图。

\begin{figure*}
  \centering
  \includegraphics[width=\textwidth]{10_f1_by_category.png}
  \caption{三模型各类别 F1 对比}
  \label{fig:10_f1_by_category}
\end{figure*}

各类别 F1 呈现严重的样本量-性能耦合特征。DoS 与 Normal 类的 F1 均在 0.77-0.88 区间，模型对样本充足且特征分布清晰的类别识别精度较高；Probe 类 F1 处于 0.65-0.74 区间，略低于前两类但仍可用；
R2L 与 U2R 类则呈现明显的两极分化：MLP 在 R2L 和 U2R 上的 F1 分别达到 0.58 和 0.56，
而 DT 和 RF 在这两类上几乎完全失效（R2L F1 仅 0.005-0.025，U2R F1 仅 0.06-0.13）。
这一巨大差距并非源于 MLP"魔法"般的判别力，而是反映了不同模型族对训练样本量的容忍度：
树模型在小样本（U2R 52 条、R2L 995 条）上几乎学不到稳定的分裂规则，
而 MLP 通过分布式表示和特征交互，能够在更少样本下提取出相对稳定的模式。
这一结果说明，\textbf{在极端类不平衡场景下，深度学习的非线性特征交互能力比树模型具有更显著的优势}。

\subsection{特征重要度收敛分析}

特征工程阶段已通过随机森林重要度筛选出 20 维核心特征，本节进一步考察 DT 与 RF 在训练完成后
给出的特征重要度排序是否一致，以判断特征选择是否具有模型无关的稳定性。
图17 并列展示两棵树模型的前 15 维特征重要度。

\begin{figure*}
  \centering
  \includegraphics[width=\textwidth]{12_feature_importance_comparison.png}
  \caption{DT 与 RF 特征重要度 Top-15 对比}
  \label{fig:12_feature_importance_comparison}
\end{figure*}

两模型在 Top-15 特征上呈现高度一致的排序：\texttt{flag\_SF}（连接状态标志）双双位列第一，
\texttt{src\_bytes}、\texttt{dst\_bytes}、\texttt{same\_srv\_rate}、\texttt{dst\_host\_srv\_count} 等流量统计特征紧随其后，
仅在第 8–15 位的次序上存在轻微差异。这一收敛现象说明：
基于信息增益或基尼不纯度的重要性度量在两个不同的树模型族上具有稳定的语义，
所选特征对各类攻击均具有判别力。换言之，特征工程的结论不依赖于具体模型的选择，
为后续模型替换与集成提供了稳健的输入基础。

\subsection{深度学习与传统机器学习}

将两类方法置于统一的测试集评估后，可以更客观地审视"深度学习必然优于传统方法"这一直觉。

在统一测试集上，DT 调优后 F1 取得 0.7621 的最高值，MLP 调优后 F1 为 0.7202，RF 为 0.6882。
深度学习方法并未在二分类任务上展示出对传统方法的系统性优势：DT 的 F1 反而高于 MLP-tuned 约 4 个百分点，CNN/LSTM 与 MLP 处于相近水平（详见表 5）。
从 AUC 指标看，RF（0.9150）> MLP-tuned（0.9074）> DT（0.7821），集成方法在概率排序层面具有判别力优势，但 DT 在准确率和 F1 上反而领先。这一对比表明，
在当前数据集上，\textbf{二分类异常检测的最优选择并非"越复杂越好"}，而是要在阈值敏感性与概率排序能力之间做权衡。

从模型架构匹配的角度，NSL-KDD 每条样本代表一次独立的网络连接，样本之间无时序或空间局部结构。
这意味着 MLP/CNN/LSTM 的归纳偏置（深度非线性、局部性、时序依赖）相对于结构扁平的表格特征均无发挥空间，
所有模型在测试集上只能依赖特征空间的相似性度量，差异仅来自函数族本身的归纳偏置。
对决策阈值为 0.5 的硬分类场景，树模型的硬决策特性反而与 NSL-KDD 的稀疏决策边界更契合。

表6 进一步比较了三模型在 5 大类多分类任务上的表现。
\textbf{MLP 的准确率达到 81.0\%，DT 与 RF 分别为 74.3\% 和 73.7\%，MLP 在 Accuracy、F1 Macro 以及 R2L/U2R 各类别 F1 上全面领先}。
这一优势在 5 大类设定下尤为突出，原因是 5 类任务需要模型对特征间的高阶交互做联合判别
（区分 R2L 与 Normal、U2R 与 DoS 等相似类别），而这正是神经网络分布式表示的强项。
树模型在分裂规则上倾向于"主特征优先"，对多特征交叉的模式捕捉能力较弱；
MLP 通过多层非线性变换，能够在小样本（R2L 995 条、U2R 52 条）下学到相对稳定的判别模式，
R2L F1（0.58）和 U2R F1（0.56）显著优于 DT/RF（R2L 0.005-0.025，U2R 0.06-0.13）。
因此，\textbf{在多分类攻击类型识别任务上，神经网络具有明显且可解释的泛化优势}。

\begin{table}[htbp]
  \centering
\caption{三模型五分类统一测试集指标对比}
  \label{tab:multiclass-comparison}
  \small
  \begin{tabular}{lcccc}
    \toprule
    模型 & Accuracy & F1 Macro & R2L F1 & U2R F1 \\
    \midrule
    MLP & \textbf{0.810} & \textbf{0.718} & \textbf{0.575} & \textbf{0.56} \\
    DT  & 0.743 & 0.488 & 0.256 & 0.08 \\
    RF  & 0.737 & 0.476 & 0.231 & 0.06 \\
    \bottomrule
  \end{tabular}
\end{table}

\subsection{综合结论}

综合上述分析，本研究在 NSL-KDD 数据集上得到以下三点主要结论。
其一，二分类场景下，DT 与 RF 是工程实践中最具竞争力的选择：
前者提供最高的 F1 与最快的推理（0.02 秒/批次），后者提供最优的概率排序与 AUC；
若以单阈值部署，DT 更优；若以概率评分支撑下游决策，RF 更优。
其二，本数据集大部分模型在二分类 hold-out 验证集上的 F1 达到 0.95 以上（MLP 0.989、CNN 0.962、LSTM 0.984，详见表 4），RF 在 5 折 GridSearchCV 上的最佳分数 0.99；仅 DT 的 GridSearchCV 最佳交叉验证分数 0.77 偏低（与 DT 树模型方差较大、容易过拟合训练子集有关）。但在统一测试集上，所有模型 F1 跌至 0.69-0.76 区间（DT 0.7621、MLP 0.7202、CNN 0.7452、LSTM 0.7161、RF 0.6882）。
这一普遍性下降源于测试集包含 15 种训练阶段从未出现的新型攻击，导致所有依赖封闭训练分布的模型都在 OOD 场景下严重贬值。
因此，对 NSL-KDD 类任务的实验结论必须在统一测试集上验证，单看训练阶段指标无法反映真实泛化能力。
其三，多分类任务在 5 类别设定下整体表现良好（DT 74.3\%、MLP 81.0\%），
但 R2L 和 U2R 因训练样本极度稀缺，F1 仍显著低于 DoS 和 Normal，
类别不平衡仍是制约小类识别能力的核心瓶颈。
此外，对比 MLP（测试集 F1=0.7202）、CNN（0.7452）、LSTM（0.7161）三种复杂度递增的神经网络，
三者性能差异极小（极差 0.029），远小于训练阶段指标暗示的"复杂度→性能"梯度。
这印证了一个重要原则：\textbf{模型不仅需要复杂，还需要与数据内在结构契合}。
NSL-KDD 表格特征之间无局部或时序依赖，CNN 的平移不变性和 LSTM 的循环结构均无发挥空间，
简单的前馈架构（MLP）反而是匹配此类扁平表格数据的最优神经网络选择。
这一发现提示我们，\textbf{盲目追求模型复杂度是工业部署中的常见误区}，应优先分析数据特性再选择架构。

基于上述结论，下一章将系统总结本研究的全部工作、归纳方法论层面的经验教训，
并对未来研究方向（如半监督学习、图神经网络、在线学习）提出展望。
